\documentclass{situationist-preprint}

\title{The Makropulos Case Reconsidered}
\msid{EXI-03-a}
\msfamily{EXI-03}
\msdiscipline{Meaning, Death and the Absurd}
\mstype{Principal manuscript}
\mskeywords{immortality; Makropulos case; Williams; categorical desires; boredom; character; psychological connectedness; meaning in life}

\begin{document}
\maketitle

\begin{abstract}
Would immortality be desirable? Bernard Williams argued that it would not: an endless life would either preserve the character of the one who lives it, in which case it would eventually become boring, or change that character so far that the future life would not be hers in the sense that matters. Critics have replied that repeatable pleasures and open-ended pursuits could sustain an endless life, and that the second horn rests on a demanding view of what survival requires. We defend Williams's conclusion in a reconstructed form. Categorical desires, those that give a person reason to go on living, are desires whose objects are projects she takes to be worthwhile; given a character that remains connected to her present one, the repertoire of such projects is bounded, and each project, pursued long enough, is either completed or reduced to repetition. Using the result of PID-06 that connectedness carries most of what matters in survival, we show that an endless life either exhausts the categorical desires that give it its point or ceases to be, in its later reaches, a life that matters to the person now. Neither horn makes death good.
\end{abstract}

\section{The problem}

In \v{C}apek's play \emph{The Makropulos Affair}, Elina Makropulos has taken an elixir that keeps her at the age of forty-two; at three hundred and forty-two, she refuses to take it again, because her life has become one of boredom and coldness. \citet{Williams1973} used the case to argue that an endless life would be meaningless, and that it is not a misfortune that we are mortal, though death is often a misfortune when it comes. His argument turns on the notion of a \emph{categorical desire}: a desire that is not conditional on one's being alive, and which can therefore give one reason to go on living, as opposed to the desire to eat when hungry, which one has only on the assumption that one will be alive to be hungry. Williams argued that an endless life would face a dilemma. If the person's character remained the same, her categorical desires would eventually be exhausted, and her life would become boring. If it changed sufficiently to renew them, the future person would not be one whose life she now has reason to want, since what she wants is that \emph{she}, with her present character, go on.

The argument has drawn three main lines of criticism. \citet{Fischer1994} argued that many pleasures are repeatable, such as those of good food, music and friendship, and that a life could be sustained indefinitely by a suitable mix of them. \citet{Chappell2009} and \citet{BortolottiNagasawa2009} argued that some pursuits, intellectual inquiry above all, are open-ended, and that boredom is a contingent psychological response that a well-constituted immortal need not have. And on the second horn, critics have pointed out that a person can survive considerable change of character, so that Williams's demand that the future person share the present character is too strong. Others, including \citet{Wisnewski2005} and \citet{Smuts2011}, have defended versions of Williams's position, and \citet{Scheffler2013} has argued that the value of many human activities depends on the scarcity of time.

We defend Williams's conclusion with a reconstructed argument. The reconstruction has two features. It identifies the categorical desires that give life its point with desires for projects the person takes to be worthwhile, in line with the account of meaning in EXI-01; and it replaces Williams's appeal to identity of character with the account of what matters in survival defended in PID-06.

\section{Framework}

\begin{definition}[Categorical desire]\label{d:cat}
A desire is \emph{categorical} if it is not conditional on the person's being alive, so that its satisfaction requires her continued life and it gives her a reason to go on living. A categorical desire \emph{gives life its point} if its object is a project that the person takes to be worthwhile for reasons independent of her desiring it.
\end{definition}

Definition~\ref{d:cat} makes explicit a restriction that Williams left implicit. A desire for another glass of wine tomorrow may be categorical in the formal sense, since one may want to be alive tomorrow in order to drink it, but it does not give life its point. EXI-01 argues that a life is meaningful to the extent that its subject engages with projects of objective worth, and EXI-17 that such engagement, not achievement, is what meaning requires. The desires that give life its point are desires to engage in such projects.

\begin{definition}[Endless life]\label{d:endless}
A life is \emph{endless} if for every time at which the person lives there is a later time at which she lives.
\end{definition}

\begin{assumption}[Bounded repertoire]\label{a:bounded}
For any character $C$, the set of project types that a person of character $C$ can engage with, in the sense of finding them worthwhile and being moved by them, is finite. A character supports engagement with a project type only by having interests, capacities and values to which that type answers.
\end{assumption}

\begin{assumption}[Exhaustion]\label{a:exhaust}
For each project type, a person who engages with it for a sufficiently long time either completes it, or comes to pursue it only as repetition of what she has already done. A project type that has been completed or reduced to repetition no longer grounds a categorical desire that gives life its point for that person.
\end{assumption}

\begin{assumption}[Connectedness]\label{a:conn}
What matters to a person, in her concern for a future life, is carried mostly by psychological connectedness between her present self and the future person, including connections of character and value (PID-06). There is a degree of connectedness below which a future life gives the person now little of what matters.
\end{assumption}

Assumption~\ref{a:bounded} is a claim about characters, not about the world. The world may contain infinitely many worthwhile things; the question is how many kinds of them a given character can care about. A person's character is a structure of interests, capacities and values of finite complexity, and the project types that answer to it are those that engage that structure. Assumption~\ref{a:exhaust} is the core of Williams's observation about Elina Makropulos. It does not say that every project can be completed; it says that if a project is not completed, then a person who pursues it long enough comes to repeat herself, and that the desire to repeat what one has already done is not a desire that gives life its point. We defend both assumptions in Section~\ref{s:obj}.

\section{The argument}

\begin{lemma}[First horn]\label{l:first}
In an endless life in which the person's character remains, at all times, connected above the threshold of Assumption~\ref{a:conn} to her present character, there is a time after which she has no categorical desire that gives life its point.
\end{lemma}

\begin{proof}
If her character remains connected above the threshold, its variation is confined within a range of characters similar to the present one, and the project types answering to characters in that range form a finite set by Assumption~\ref{a:bounded}, since the range is the range of a bounded variation on a character of finite complexity. Each of these types, pursued for a sufficiently long time, is completed or reduced to repetition by Assumption~\ref{a:exhaust}. The life is endless (Definition~\ref{d:endless}), so there is time for each type in the finite set to be pursued for as long as Assumption~\ref{a:exhaust} requires, and after the last of these periods no type in the set grounds a categorical desire that gives life its point. Projects of other types do not engage her, by Assumption~\ref{a:bounded}. So there is a time after which she has none.
\end{proof}

\begin{lemma}[Second horn]\label{l:second}
In an endless life in which the person's character does not remain connected above the threshold, there are times after which the life gives her, as she now is, little of what matters.
\end{lemma}

\begin{proof}
If her character falls at some time below the threshold of connectedness to her present character, then by Assumption~\ref{a:conn} the life after that time gives her now little of what matters. The future person may be psychologically continuous with her, through overlapping chains of connections, and may even be identical with her, since PID-01 makes continuity, not connectedness, the criterion of identity. But PID-06 shows that identity and continuity without connectedness carry little of what matters. So the later reaches of the life are not something she now has reason to want for her own sake, though they may be good for the person who lives them.
\end{proof}

\begin{theorem}[Immortality is not desirable]\label{t:main}
An endless life is not desirable for the one who would live it: it would either eventually exhaust the categorical desires that give life its point, or in its later reaches cease to be a life that gives her what matters.
\end{theorem}

\begin{proof}
Every endless life falls under Lemma~\ref{l:first} or Lemma~\ref{l:second}, according as the person's character remains connected above the threshold at all times or does not. Under the first, it contains an endless period without categorical desires that give life its point, so the person's life becomes, from some time on, pointless; under the second, its later reaches give her now little of what matters. In neither case does the endless life offer her what she wants when she wants not to die: the continuation of a life with point, for herself as she is.
\end{proof}

Theorem~\ref{t:main} is a claim about an endless life, not about a longer one. It is consistent with the conclusion of EXI-02 that death is bad for the one who dies when it deprives her of goods she would otherwise have had. A person who dies while she still has categorical desires that give her life its point is deprived by her death; a person whose categorical desires are exhausted may not be. Williams held both claims, and the reconstruction shows why they do not conflict: the deprivation of EXI-02 diminishes as the life approaches the point at which Lemma~\ref{l:first} applies, and an endless life is guaranteed to reach that point or to lose connectedness first.

\section{Objections and replies}\label{s:obj}

\begin{objection}[Repeatable pleasures]\label{o:fischer}
Many goods are repeatable without loss: one does not tire of good food, music or the company of friends. A life could be sustained endlessly by these \citep{Fischer1994}.
\end{objection}

\begin{reply}
Repeatable pleasures are goods, and an endless life containing them would contain goods at every time. But the desire for them does not give life its point in the sense of Definition~\ref{d:cat}: the person who wants to go on living in order to have another good meal wants to go on living, and wants pleasure while she does, but the meal is not a project she takes to be worthwhile independently of her wanting it. A life sustained only by such desires is the life Elina Makropulos describes, in which things still happen and some of them are pleasant. Fischer's examples of friendship are different, and fall under the next objection: friendship is a project, and its renewal depends on the friends' changing.
\end{reply}

\begin{objection}[Open-ended pursuits]\label{o:open}
Some pursuits never run out. Mathematics contains infinitely many theorems, and there is always more to understand about the world and about other people. A person engaged in inquiry would never be reduced to repetition \citep{Chappell2009,BortolottiNagasawa2009}.
\end{objection}

\begin{reply}
The objects of inquiry are inexhaustible, but Assumption~\ref{a:exhaust} concerns the person's engagement. A mathematician with a fixed character works in some areas, with some methods, moved by some kinds of problem. Over an endless life the new theorems available to her become, in the respects that engage her, instances of kinds she has already proved; to find them new she would have to acquire new interests and capacities, and a person who has acquired enough of these has changed her character in the way that leads to Lemma~\ref{l:second}. The objection shows that the world does not run out of things worth doing. It does not show that a connected character does not run out of ways to care about them.
\end{reply}

\begin{objection}[Boredom is contingent]\label{o:bored}
The argument relies on boredom, a psychological response that an immortal could lack, by constitution or by a suitably managed life.
\end{objection}

\begin{reply}
Boredom is the symptom, not the premise. An immortal who did not feel bored after her categorical desires were exhausted would not thereby have them again; she would have a life without point that does not trouble her. Theorem~\ref{t:main} concerns whether the endless life is desirable, and a life whose pointlessness is not felt is not thereby a life with a point. EXI-10 argues that boredom signals a mismatch between capacities and activities; an immortal without boredom would lack the signal, not the mismatch.
\end{reply}

\begin{objection}[Gradual change]\label{o:grad}
Character can change gradually and indefinitely while each stage is strongly connected to the next. Such a life would always have categorical desires, and at each time its subject would be glad to go on.
\end{objection}

\begin{reply}
This is the second horn. At each time the person is glad to go on, and each stage of the life may be good for the person who lives it. But the question is whether the endless life is desirable for the person who now considers it. The person a thousand years on is connected to her now only through a chain, and by PID-06 continuity through a chain carries little of what matters. The life described is something like a succession of lives, each of which has reason to want the next to begin; that is not what someone wants who wants not to die.
\end{reply}

\section{What the result does not establish}

The argument does not show that any finite life is too long, or that there is a determinate age at which categorical desires give out; Assumption~\ref{a:exhaust} speaks of a sufficiently long time, and for most people that time may be very long. It does not show that death is good, or that it is not bad for the one who dies (EXI-02). It does not apply to forms of afterlife that do not consist in an endless succession of experiences and projects, such as the timeless states of some theological traditions. And it does not show that the limits of a mortal life are what give it its shape, which is a separate claim discussed in EXI-19.

\section{Conclusion}

An endless life would eventually exhaust the categorical desires that give life its point, or else cease, in its later reaches, to be a life that gives its subject now what matters to her. Williams's dilemma survives its critics once categorical desires are understood as desires for worthwhile projects, and once his appeal to sameness of character is replaced by the claim that connectedness carries most of what matters. Immortality is not desirable, though death, when it comes, usually takes something from the one who dies.

\begin{thebibliography}{99}\small
\bibitem[Bortolotti and Nagasawa(2009)]{BortolottiNagasawa2009} Bortolotti, L. and Nagasawa, Y. (2009). Immortality without boredom. \emph{Ratio}, 22(3), 261--277.
\bibitem[Chappell(2009)]{Chappell2009} Chappell, T. (2009). Infinity goes up on trial: Must immortality be meaningless? \emph{European Journal of Philosophy}, 17(1), 30--44.
\bibitem[Fischer(1994)]{Fischer1994} Fischer, J.~M. (1994). Why immortality is not so bad. \emph{International Journal of Philosophical Studies}, 2(2), 257--270.
\bibitem[Scheffler(2013)]{Scheffler2013} Scheffler, S. (2013). \emph{Death and the Afterlife}. Oxford University Press.
\bibitem[Smuts(2011)]{Smuts2011} Smuts, A. (2011). Immortality and significance. \emph{Philosophy and Literature}, 35(1), 134--149.
\bibitem[Williams(1973)]{Williams1973} Williams, B. (1973). The Makropulos case: Reflections on the tedium of immortality. In \emph{Problems of the Self}, 82--100. Cambridge University Press.
\bibitem[Wisnewski(2005)]{Wisnewski2005} Wisnewski, J.~J. (2005). Is the immortal life worth living? \emph{International Journal for Philosophy of Religion}, 58(1), 27--36.
\end{thebibliography}

\end{document}
